\section{Otros tipos de Ciclo de Vida del Análisis de Datos}
\subsection{EMC/Dell EMC - Proceso cíclico}
Descubrimiento - Preprocesamiento de datos - Planificación de modelos - Construcción de Modelos - Comunicación (de resultados) - Puesta en funcionamiento
Las fases no son hitos estáticos = dan paso al siguiente y luego se repiten

\subsection{SAS - Iterativo}
Preguntar - Preparar - Explorar - Modelar - Implementar - Actuar - Evaluar [evaluar soluciones para posteriormente volver a preguntar]
Producir resultados repetibles, fiables y predictivos

\subsection{Basado en Proyectos}
Identificación del problema - Diseño de requisito de datos - Preprocesamiento de datos - Realización del análisis de datos - Visualización de datos 

\subsection{Ciclo del Big Data}
Evaluación del caso de negocio - Identificación de datos - Adquisición y Filtrado de Datos - Extracción de datos - Validación y Limpieza de datos - Agregación y representación de datos - Análisis de datos - Visualización de datos - Utilización de resultados del análisis
\newpage

\paragraph{crash course en notación matemática}

Tipos de número y constantes importantes: Naturales \N, Racionales \Q, Reales \R, Complejos \C, número imaginario \textit{\textbf{i}}, velocidad de la luz \textbf{\textit{c}}, número de euler \textbf{e} (que se puede escribir solo o en forma de función tipo \textbf{\textit{f(}\(e^{st}\)}))

 
Teoría de conjuntos: Existe \(\exists\), Existe un único elemento \(\exists!\), No existe \(\nexists\), Pertenece \(\in\), No pertenece \(\notin\), Para todo en un conjunto \(\forall\), y lógico \(\land\), o lógico \(\lor\), conjunto vacío \(\emptyset\)

Letras griegas que se usan harto (en minúscula): generalmente se ocupa \(\alpha, \beta, \theta, \phi\) para ángulos, cantidades pequeñas con épsilon \(\epsilon,  \varepsilon\), diferencia con delta en may/min \(\Delta / \delta\), coeficiente de roce mu \(\mu\), tautología aka verdad absoluta tau \(\tau\), longitud (también coeficiente random) lambda \(\lambda\), desviación estándar sigma \(\sigma\)

Vectores: letras con flechitas arriba, los más comunes en física son \(\vec{a}\)celeración, \(\vec{v}\)elocidad, \(\vec{r}\)adio/posición. Los vectores unitarios tienen gorrito, los más comunes son \(\hat{\imath}, \hat{\jmath}\, \hat{k}, \hat{\phi}, \hat{\theta}\)

Operadores: derivada parcial \(\partial\), operador diferencial nabla \(\nabla\), integral \(\int\), sumatoria con sigma mayúscula \(\Sigma^{n}_{i=0}\), unión de conjuntos \(\bigcup^{n}_{i=1}\), producto punto \(\cdot\)

Funciones: generalmente cualquier letra(y algo encerrado en paréntesis) tipo f(x), g(x), r(t), $\omega (\theta)$

Aquí algunas ideas para los nombres
\begin{itemize}
    \item {Ál\textit{\textbf{g}}(\(\varepsilon \beta r \alpha\)) \textbf{\textit{f(}\(e^{st}\)})}
    \item {\R ichar\(\partial\) Parque\(\vec{r}\)}
    \item {\C hip\textbf{e} Li\(\beta\)r\textbf{e}}
    \item {\Q ui\(\forall i \hat{k}\) o \Q ui\(\lor\)\textbf{\textit{i}}k}
    \item {\N isu\(y^{a\textit{\textbf{i}}}\) o \N isuy(\(\alpha\)\textit{\textbf{i}})}
    \item {Man\(\hat{j}\)u o \(\mu \alpha\)nj\(\lor\)}
    \item {Yer\(\hat{k}\)o \(\bigcup^{rr}_a\) o Yerk\(\emptyset\) \(\bigcup rr\alpha\)}
    \item {\(\mu a \tau ar\)  \(\alpha\) g(rax)}
    \item {\(\int\)a \textbf{e}strat\(\textbf{e}^{g\textit{i}a}\) \(\delta \epsilon\)l cara\textbf{\textit{c}}ol}
    \item {Trasa\(\tau l \vec{a}\)nt\textit{\textbf{i}}ca o Trasa\(\tau l\vec{a}\)nti\textit{\textbf{c}}a}
    \item {Tri\(\alpha ng(ul\theta)\) \(\delta \varepsilon\) P\(\textbf{e}^{n}\) \(\cdot\) rose}
\end{itemize}












\newpage